\section{Methodology and theory}
\label{sec:methodology}

\subsection{Protocol definitions}

Let $\{\phi_t^{(n)}\}_{n=1}^{\Nshots}$ denote the shot-level readout features at time step $t$.

\paragraph*{EV.}
All shots are averaged once:
\begin{equation}
x_t^{\mathrm{EV}} = \frac{1}{\Nshots}\sum_{n=1}^{\Nshots} \phi_t^{(n)}.
\end{equation}
Each time step contributes one training row.

\paragraph*{Raw.}
Every shot is kept separately:
\begin{equation}
x_{t,n}^{\mathrm{Raw}} = \phi_t^{(n)}.
\end{equation}
Each time step contributes $\Nshots$ training rows, but each row carries the full shot noise.

\paragraph*{\Split{}.}
Choose a divisor $k$ of $\Nshots$, partition the shots into $G=\Nshots/k$ disjoint groups, and form
\begin{equation}
x_{t,g}^{(k)} = \frac{1}{k}\sum_{n\in G_g}\phi_t^{(n)}, \qquad g=1,\dots,G.
\label{eq:splitdef}
\end{equation}
Each grouped row is paired with the same target $y_t$ during training. At validation and test time, the grouped predictions are averaged back to one prediction per labeled time step.

The group size $k$ is selected by chronological validation over the proper divisors of $\Nshots$. This keeps the search discrete, transparent, and budget fair.

\subsection{Two additional controls}

We use two stronger controls beyond the three protocol endpoints.

\paragraph*{EV-dup.}
EV-dup duplicates the same EV row enough times to match the grouped row count. This isolates ``more rows'' from ``more distinct grouped views.'' If the gain came only from the nominal row count, EV-dup should match \Split{}.

\paragraph*{Variance-aware baselines.}
We also include EV-NA and Split-NA, diagonal noise-aware ridge baselines of the form
\begin{equation}
\|y-Xw\|_2^2 + \lambda \|w\|_2^2 + \gamma\, w^\top \widehat{\Sigma}_{\epsilon,k} w,
\label{eq:noiseaware}
\end{equation}
with $\gamma$ chosen by validation. These baselines are not meant as a final word on finite-shot mitigation, but they test whether the grouped-view gain can be explained away by a simple variance-aware penalty alone.

\subsection{Grouped-view surrogate}

To understand when grouping should help, we use a grouped-view surrogate adapted from classical measurement-error models \cite{fuller1987measurement,carroll2006measurement}. Let
\begin{equation}
y_t = \beta_\star^\top z_t + \xi_t,
\label{eq:latent}
\end{equation}
where $z_t\in\mathbb{R}^d$ is a latent clean feature vector and $\xi_t$ is residual target noise. A grouped average over $k$ shots is modeled as
\begin{equation}
x_{t,g}^{(k)} = z_t + \varepsilon_{t,g}^{(k)}, \qquad
\mathbb{E}\varepsilon_{t,g}^{(k)}=0, \qquad
\mathrm{Cov}(\varepsilon_{t,g}^{(k)})=\frac{\Sigma_\epsilon}{k}.
\label{eq:groupedview}
\end{equation}
EV corresponds to $k=\Nshots$, Raw to $k=1$, and \Split{} to $1<k<\Nshots$.

The point of this model is narrow and deliberate. It is not a claim that all QRC feature distributions are isotropic or independent. It is a mechanism model whose job is to answer one question: what exact transformation does grouping apply to the geometry of the readout problem, and how does that interact with feature noise?

\subsection{Exact aspect-ratio lift}

\begin{proposition}[Exact grouped-view aspect-ratio lift]
Let $\rhoEV=\Ttrain/d$ be the EV aspect ratio, where $\Ttrain$ is the number of labeled training time steps before grouping and $d$ is the readout-feature dimension. If a grouped protocol uses group size $k$ and therefore contributes $G=\Nshots/k$ rows per time step, then the grouped design has aspect ratio
\begin{equation}
\rhoK = \frac{G\Ttrain}{d} = \frac{\Nshots}{k}\rhoEV.
\label{eq:rhok}
\end{equation}
\end{proposition}

\begin{proof}
Each of the $\Ttrain$ training time steps contributes $G=\Nshots/k$ grouped rows. The total grouped row count is therefore $G\Ttrain$, while the feature dimension remains $d$. Dividing by $d$ yields Eq.~\eqref{eq:rhok}.
\end{proof}

This proposition is exact. It says that grouping moves the readout problem to the right on the aspect-ratio axis.

\subsection{Attenuation and row-count tradeoff}

The exact lift above is useful but, by itself, too weak to guide the choice of $k$. To sharpen the theory, we add the standard attenuation mechanism from noisy-covariate regression. In the scalar case, let $z=x+\epsilon$ with
\begin{equation}
\mathrm{Var}(\epsilon)=\frac{\sigma_\epsilon^2}{k},
\end{equation}
and let $y=w_\star x+\nu$. The population regression slope of $y$ on $z$ is
\begin{equation}
w_\infty(k)=\beta_0(k)w_\star,
\qquad
\beta_0(k)=\frac{\sigma_x^2}{\sigma_x^2+\sigma_\epsilon^2/k}.
\label{eq:beta0}
\end{equation}
Writing $a:=\sigma_\epsilon^2/\sigma_x^2=1/\mathrm{SNR}_m$, this becomes
\begin{equation}
\beta_0(k)=\frac{k}{k+a}.
\label{eq:beta0a}
\end{equation}
Small $k$ therefore increases the effective row count but also increases attenuation.

The row-count side is exact. Because EV contributes $\rhoEV d$ fully averaged rows and grouping yields $\Nshots/k$ grouped rows per time step, the grouped effective row count is
\begin{equation}
N_{\mathrm{eff}}(k)=\frac{\rhoEV d\,\Nshots}{k}.
\label{eq:neff}
\end{equation}
Following the standard surrogate style in the measurement-error and random-design ridge literature \cite{hsu2012random,dobriban2018high}, we combine the two effects in a simple proxy:
\begin{equation}
\mathcal{E}(k) \approx c_{\mathrm{att}}\left(\frac{a}{k+a}\right)^2 + c_{\mathrm{rows}}\frac{k}{\rhoEV\Nshots}.
\label{eq:proxy}
\end{equation}
The first term penalizes attenuation; the second penalizes too few grouped rows.

\begin{proposition}[Warm-start proxy for the preferred group size]
Under the proxy in Eq.~\eqref{eq:proxy}, the continuous minimizer is
\begin{equation}
k_\star^{\mathrm{cont}} = \left(\frac{2c_{\mathrm{att}}}{c_{\mathrm{rows}}}a^2\rhoEV\Nshots\right)^{1/3} - a.
\label{eq:kstar}
\end{equation}
In practice we clip this value to $[1,\Nshots]$ and round to the nearest divisor of $\Nshots$ before validation.
\end{proposition}

\begin{proof}
Differentiating Eq.~\eqref{eq:proxy} and setting the derivative to zero yields
\begin{equation}
(k+a)^3 = \frac{2c_{\mathrm{att}}a^2\rhoEV\Nshots}{c_{\mathrm{rows}}},
\end{equation}
from which Eq.~\eqref{eq:kstar} follows immediately.
\end{proof}

This proposition is still a surrogate result, not a universal optimum theorem. Its value is practical. It turns the qualitative statement ``grouping helps when EV is row-limited'' into a usable recipe for initializing the discrete search over divisors of $\Nshots$.

\subsection{Conditioning landmarks}

For isotropic random-design asymptotics, the lower edge of the sample covariance spectrum is
\begin{equation}
\gMP(\rho)=\left(1-\rho^{-1/2}\right)_+^2.
\label{eq:mp}
\end{equation}
Because $\rhoK>\rhoEV$ whenever $k<\Nshots$, grouping strictly lifts the MP lower edge. We use $\rho=1$ and $\rho=4$ as convenient conditioning landmarks: $\rho=1$ is the EV rank threshold, while $\rho=4$ corresponds to $\gMP(4)=0.25$. The theory therefore predicts a simple regime picture: grouping should matter most when EV is still conditioning-limited and should matter less once EV is comfortably above the edge.

What the theory \emph{does not} do is determine a universal numeric test-error crossover. The exact lift, the attenuation factor, and the warm-start proxy are the theoretically useful parts. Whether the final test error is lower remains empirical and task dependent, which is exactly what the next sections test.
